\chapter*{Phụ lục}
\markboth{PHỤ LỤC}{PHỤ LỤC}
\addcontentsline{toc}{chapter}{Phụ lục}

\section*{Phụ lục A --- DDL các bảng}

\begin{sloppypar}Xem file \texttt{sql/01\_\allowbreak{}schema\_\allowbreak{}audit.sql} (bảng \texttt{audit\_logs} partitioned, \texttt{security\_alerts}) và \texttt{sql/02\_\allowbreak{}schema\_\allowbreak{}business.sql} (bảng \texttt{orders}, \texttt{products}).\end{sloppypar}

\textbf{Tóm tắt schema:}

\begin{center}
\captionof{listing}{DDL --- bảng audit logs và security alerts}
\label{lst:appendix-ddl}
\begin{Verbatim}[frame=single, fontsize=\scriptsize, numbers=left, numbersep=5pt, breaklines]
-- Bảng audit chính (partitioned)
CREATE TABLE audit_logs (
    id          BIGSERIAL,
    table_name  TEXT      NOT NULL,
    operation   TEXT      NOT NULL,   -- INSERT / UPDATE / DELETE
    user_name   TEXT,
    old_data    JSONB,
    new_data    JSONB,
    changed_at  TIMESTAMP NOT NULL DEFAULT CURRENT_TIMESTAMP,
    prev_hash   BYTEA,                -- tamper-evident chain
    hash        BYTEA,
    PRIMARY KEY (id, changed_at)
) PARTITION BY RANGE (changed_at);

-- Partition theo tháng (ví dụ tháng 4/2026)
CREATE TABLE audit_logs_2026_04
PARTITION OF audit_logs
FOR VALUES FROM ('2026-04-01') TO ('2026-05-01');

-- Bảng security_alerts (autonomous commit via dblink)
-- Xem sql/01_schema_audit.sql

-- Bảng audit_ddl_logs (DDL audit qua event trigger --- sql/09_audit_ddl.sql)
CREATE TABLE public.audit_ddl_logs (
    id          BIGSERIAL PRIMARY KEY,
    command_tag TEXT      NOT NULL,   -- CREATE TABLE, ALTER TABLE, DROP TABLE, ...
    object_type TEXT,                 -- TABLE / INDEX / FUNCTION
    object_name TEXT,                 -- full object name within schema
    command_sql TEXT,                 -- command_tag || ' ' || object_identity
    executed_at TIMESTAMP NOT NULL DEFAULT CURRENT_TIMESTAMP,
    user_name   TEXT      NOT NULL DEFAULT session_user
);
\end{Verbatim}
\end{center}

\section*{Phụ lục B --- Code PL/pgSQL}

Xem các file SQL trong thư mục \texttt{sql/}:

\begin{center}
\captionof{table}{Code PL/pgSQL chính}
\label{tab:code-files}
\begin{tabularx}{\textwidth}{l>{\raggedright\arraybackslash}X}
\toprule
\textbf{File} & \textbf{Nội dung} \\
\midrule
\texttt{sql/04\_audit\_function.sql} & \texttt{func\_audit\_trigger()} --- generic audit trigger, SECURITY DEFINER, \texttt{session\_user} \\
\texttt{sql/05\_security\_immutability.sql} & \texttt{func\_prevent\_audit\_change()} --- immutability trigger và dblink autonomous alert \\
\texttt{sql/06\_hash\_chain.sql} & \texttt{func\_audit\_hash\_chain()} --- SHA-256 hash chain; \texttt{func\_verify\_hash\_chain()} --- verifier \\
\texttt{sql/09\_audit\_ddl.sql} & \texttt{func\_audit\_ddl()} --- event trigger audit DDL (CREATE/ALTER), lưu vào \texttt{audit\_ddl\_logs} \\
\bottomrule
\end{tabularx}
\end{center}

\section*{Phụ lục C --- Script benchmark và verify}

\begin{center}
\captionof{table}{Script benchmark và kiểm thử}
\label{tab:benchmark-scripts}
\begin{tabularx}{\textwidth}{l>{\raggedright\arraybackslash}X}
\toprule
\textbf{File} & \textbf{Mô tả} \\
\midrule
\texttt{bench/update\_orders.sql} & pgbench script: UPDATE ngẫu nhiên 1 đơn hàng \\
\texttt{bench/run\_baseline.sh} & Chạy 5 lần baseline (trigger disabled), in TPS/latency \\
\texttt{bench/run\_proposed.sh} & Chạy 5 lần proposed (trigger enabled), tính overhead vs baseline \\
\texttt{bench/analyze\_performance.sh} & Phân tích \texttt{pg\_stat\_statements}: top queries, audit impact \\
\texttt{bench/monitor.sh} & Real-time monitoring: connections, TPS, WAL, cache hit ratio \\
\texttt{verify/q\_partition\_pruning.sql} & Đo DROP PARTITION vs DELETE; EXPLAIN partition pruning \\
\texttt{verify/q\_gin\_comparison.sql} & Đo JSONB query có/không GIN index (auto drop/recreate) \\
\texttt{verify/q\_security\_demo.sql} & Demo 3 kịch bản tấn công (PASS/FAIL) \\
\texttt{verify/q\_ddl\_audit.sql} & Demo DDL audit: lịch sử DDL, live capture, thống kê theo user/command \\
\texttt{verify/q\_jsonb\_examples.sql} & Ví dụ truy vấn JSONB phục vụ kiểm toán \\
\texttt{verify/q\_hash\_verifier.sql} & Kiểm tra tính toàn vẹn hash chain \\
\bottomrule
\end{tabularx}
\end{center}

\section*{Phụ lục D --- Truy vấn mẫu phục vụ kiểm toán}

\begin{center}
\captionof{listing}{D1. Lịch sử thay đổi 50 dòng gần nhất trên bảng orders}
\label{lst:query-d1}
\begin{Verbatim}[frame=single, fontsize=\scriptsize, numbers=left, numbersep=5pt, breaklines]
SELECT operation, user_name,
       old_data->>'status' AS old_status,
       new_data->>'status' AS new_status,
       changed_at
FROM audit_logs
WHERE table_name = 'public.orders'
ORDER BY changed_at DESC LIMIT 50;
\end{Verbatim}
\end{center}

\begin{center}
\captionof{listing}{D2. Truy vấn JSONB --- tìm đơn hàng chuyển sang PAID}
\label{lst:query-d2}
\begin{Verbatim}[frame=single, fontsize=\scriptsize, numbers=left, numbersep=5pt, breaklines]
SELECT id, user_name, old_data->>'status' AS old_status, changed_at
FROM audit_logs
WHERE table_name = 'public.orders'
  AND new_data @> '{"status": "PAID"}';
\end{Verbatim}
\end{center}

\begin{center}
\captionof{listing}{D3. Truy vết actor --- ai thay đổi nhiều nhất hôm nay}
\label{lst:query-d3}
\begin{Verbatim}[frame=single, fontsize=\scriptsize, numbers=left, numbersep=5pt, breaklines]
SELECT user_name, operation, count(*) AS cnt
FROM audit_logs
WHERE changed_at::date = current_date
GROUP BY user_name, operation
ORDER BY cnt DESC;
\end{Verbatim}
\end{center}

\begin{center}
\captionof{listing}{D4. Kiểm tra tính toàn vẹn hash chain}
\label{lst:query-d4}
\begin{Verbatim}[frame=single, fontsize=\scriptsize, numbers=left, numbersep=5pt, breaklines]
SELECT chain_ok, count(*)
FROM func_verify_hash_chain('public.orders')
GROUP BY chain_ok;
-- chain_ok = true: hợp lệ; false: có dòng bị sửa trực tiếp tầng file
\end{Verbatim}
\end{center}

\begin{center}
\captionof{listing}{D5. Phân tích DDL audit --- lịch sử thay đổi schema}
\label{lst:query-d5}
\begin{Verbatim}[frame=single, fontsize=\scriptsize, numbers= left, numbersep=5pt, breaklines]
SELECT command_tag, object_type, object_name, user_name, executed_at
FROM audit_ddl_logs
WHERE executed_at >= '2026-04-01'
ORDER BY executed_at DESC;
\end{Verbatim}
\end{center}
